%=============================================================================!
% 4.1  THE SHADOW OF A CIRCLE                                                 !
%=============================================================================!
\section{The shadow of a circle}
\label{sec:os-circle}

A block of mass $m$ on a spring of stiffness $k$ moves, by
\cref{sec:ne-spring}, as
\begin{equation}
   x(t) = A\cos(\omega t + \phi) , \qquad \omega = \sqrt{\frac{k}{m}} ,
   \label{eq:os-harmonic}
\end{equation}
where the amplitude $A$ and the phase $\phi$ are fixed by the position
and velocity at the start. This motion is called simple harmonic motion,
and the system that performs it a harmonic oscillator. This section
gives the formula a picture that makes its phase, its velocity and its
period easy to read.

\subsection*{The reference circle}

\Cref{sec:mo-circle} followed a point going round a circle at a steady
angular speed, and found that its projection on the $x$ axis moves back
and forth as a sinusoid. Now run the argument the other way. Let a point $P$ go anticlockwise
round a circle of radius $A$ centred on the origin, at the angular speed
$\omega$, starting at the angle $\phi$ from the $x$ axis, so that at time
$t$ its angle is
\begin{equation*}
   \theta(t) = \omega t + \phi .
\end{equation*}
By the definition of the cosine, scaled to a circle of radius $A$
(\cref{sec:mo-trig}), the $x$ coordinate of $P$ is $A\cos(\omega t +
\phi)$, which is \cref{eq:os-harmonic}. The block therefore moves exactly
as the shadow of $P$ on the $x$ axis, the point directly below or above
it (\cref{fig:os-circle}). The circle is called the reference circle of
the oscillation, and the phase $\phi$ is the angle at which $P$
starts.

The velocity follows the same picture. \Cref{eq:mo-circle} differentiated
$R(\cos\omega t, \sin\omega t)$; here the radius is $A$ and the angle is
$\omega t + \phi$, which also grows at the rate $\omega$, so the chain rule
gives the same factor $\omega$ and the velocity of $P$ is $A\omega\,(-\sin
\theta, \cos\theta)$. Its $x$ component, $-A\omega\sin(\omega t + \phi)$,
is the derivative of $x(t)$: the velocity of the block is the shadow of the
velocity of $P$.
The point $P$ always moves at the speed $A\omega$, tangent to the circle.
Where $P$ crosses the top or the bottom of the circle it moves parallel to
the $x$ axis, so the shadow moves at the full speed $A\omega$; this is the
centre of the swing. Where $P$ crosses the $x$ axis it moves straight up
or down, so the shadow stands still for an instant; these are the turning
points, $x = \pm A$. Likewise the acceleration of $P$ points towards the
centre with size $\omega^2 A$ (\cref{sec:mo-circle}), and its $x$
component is $-\omega^2 x$, the acceleration of the block.

\begin{figure}[htb]
\centering
\includegraphics{ch04_oscillations/circle}
\caption{The oscillation as the shadow of a circle. A block with $\omega =
\SI{10}{\radian\per\second}$ starts \SI{3}{\centi\metre} out, moving
outwards at \SI{0.4}{\metre\per\second}. (a) The point $P$, going
anticlockwise round a circle of radius $A = \SI{5}{\centi\metre}$ from
$\phi = \SI{-0.927}{\radian}$, at $t = 0$, 0.1, 0.2 and \SI{0.3}{\second}
(dots), and its shadow on the $x$ axis (squares), one colour per moment.
(b) The shadow's position against time, with time running downwards so
that each square lies directly below the square of the same colour in
panel (a).}
\label{fig:os-circle}
\end{figure}

\subsection*{An example}

Take the block of \cref{sec:ne-spring}, with $m = \SI{0.5}{\kilogram}$
and $k = \SI{50}{\newton\per\metre}$, so that $\omega =
\SI{10}{\radian\per\second}$, started \SI{3}{\centi\metre} from its
resting place and moving outwards at \SI{0.4}{\metre\per\second}. That
section found $A = \SI{5}{\centi\metre}$, $\cos\phi = x(0)/A = 0.6$ and
$\phi = \arcsin(-0.8) = \SI{-0.927}{\radian}$. On the reference circle,
$P$ starts \SI{0.927}{\radian}, or $0.927\times\ang{180}/\pi =
\ang{53.1}$, below the positive $x$ axis, where
its shadow is at $A\cos\phi = 5\times0.6 = \SI{3}{\centi\metre}$. Moving
anticlockwise, $P$ climbs towards the axis, so the shadow moves outwards,
as the block does. $P$ reaches the axis when $\omega t + \phi = 0$, at $t
= 0.927/10 = \SI{0.0927}{\second}$, and there the block turns back, at
$x = A = \SI{5}{\centi\metre}$.

$P$ goes once round the circle while $\omega t$ grows by $2\pi$, so the
period is $2\pi/\omega = 2\pi/10 = \SI{0.628}{\second}$, and the
frequency, the number of oscillations per second (\cref{sec:mo-trig}),
is
\begin{equation*}
   \nu = \frac{\omega}{2\pi} = \frac{10}{2\pi} = \SI{1.592}{\hertz} ,
\end{equation*}
written with the Greek letter nu. The angular frequency $\omega$ is the
rate at which $P$ turns, in radians per second, and is $2\pi$ times the
frequency because each oscillation is one full turn of $2\pi$ radians.

\begin{selfcheck}
A block oscillates with $\omega = \SI{10}{\radian\per\second}$ and $A =
\SI{5}{\centi\metre}$. Where is the block, and which way and how fast is
it moving, when the angle $\omega t + \phi$ of its point $P$ equals
$\pi/2$?
\end{selfcheck}
